\documentclass{article}
\usepackage{amsmath,amssymb,amsthm,hyperref}
\title{Native law specifications and shared event assembly}
\date{16 November 2026}
\begin{document}\maketitle
\paragraph{Status.} The row-adaptive and globally scoped constructions and their obstruction arguments passed independent written-proof review, including the choice-scope, loss-mask and pure-input exactness boundaries. The target is a structural bridge between joint native assembly, chronological substitution and event restriction. An event-restriction adjunction is not by itself a trace theorem; no unrestricted feedback or full v0.2 completion is claimed.
\section{The obstruction to discrete law families}
Let $\mathcal W_\le$ have measured raw carriers $(W,p)$, finitely supported nonnegative finite measures, and actual functions $r$ satisfying $r_*p\le q$. Coefficients may be trivial for this obstruction. Its empty carrier is initial. A multiversal product family must therefore have exactly one component: the unique cone from the initial object factors uniquely through each component. Thus a missing ordinary product cannot be repaired by discrete families.

Two fair-bit objects have no product. Zero-mass singleton probes force any product's raw carrier to be the four ordered pairs of bits. A mass-$1/2$ singleton cone to each pair forces each product atom to have mass at least $1/2$. A projection then has mass at least one above each bit, violating the fair marginal bound $1/2$. Hence no multiproduct exists. In particular, the exact-effect all-coupling construction cannot simply be reused with subpreserving maps.

For a fixed finite diagram with raw limit $Z$ and coefficient limits, there is nevertheless an ordered universal description. Let $K_D$ be its poset of admissible joint sublaws, and $J(t)$ the canonical cone with carrier $Z$ and law $t$. A cone $c$ has a unique raw/coefficient factor $r$ and induced law $t_c=r_*p$. Then
\[\mathrm{Cone}(c,J(t))\cong\{* : t_c\le t\}.\]
Thus $c\mapsto t_c$ is left adjoint to the full inclusion of canonical joint-law cones. This is a reflective ordered interface, not a discrete multiversal family. The diagram has a limit exactly when $K_D$ has a greatest element.
\section{Native specification fibres}
Fix a finitely complete coefficient category $\mathcal V$. Over a label set $X$, an object is
\[(W\xrightarrow{\ell}X,\ (C_w),\ K),\]
where $K$ is a nonempty downward-closed convex set of finitely supported finite nonnegative measures on $W$. Convex means closure under finite barycentric mixtures; no topological closure or probability distribution on the set $K$ is imposed. Nonemptiness and downward closure imply $0\in K$. Raw zero-weight witnesses and their coefficients are retained. A map consists of an actual $r$ over $X$ and actual coefficient arrows $C_w\to D_{r(w)}$ with $r_*K\subseteq L$.

The principal embedding $(W,p,C)\mapsto(W,\downarrow p,C)$ is fully faithful for subpreserving maps: $r_*(\downarrow p)\subseteq\downarrow q$ iff $r_*p\le q$. On normalized probability objects this recovers exact maps because both measures have total one. It does not make the arbitrary finite-measure exact doctrine locally full.

\paragraph{Finite limits.} For a finite diagram take the raw limit $Z$ in $\mathrm{Set}/X$, coefficient limits at its points, and
\[K_Z=\{t:(\pi_j)_*t\in K_j\ \text{for every }j\}.\]
This set is nonempty, lower and convex. A cone factors uniquely on raw carriers and coefficient arrows; its factor is admissible precisely because its projections are. This proves the categorical limit property. The terminal object has raw carrier $X$, terminal coefficients and all finite measures on $X$. The result retains every admissible joint sublaw; it chooses none. A selected law is additional data, expressible by its principal object mapping into the specification.
\section{Chronological bind and its forced convexity}
For chronological coefficient assembly assume additionally that $\mathcal V$ is monoidal. For an outer specification $K$ on $W$ and nonempty lower convex continuation specifications $L_w$ on $Z_w$, define bind on the retained raw sum $\sum_w Z_w$ by the laws
\[t(w,z)=p(w)q_w(z),\qquad p\in K,\quad q_w\in L_w.\]
Only finitely many positive outer weights matter; at other raw witnesses choose zero. Finite support is preserved. This set is lower: replace each row by the corresponding smaller row divided by $p(w)$ where positive. It is convex: a mixture of factorizations has outer law $p=\sum_i\lambda_i p_i$ and row law
\[q_w=\sum_i\frac{\lambda_i p_i(w)}{p(w)}q_{i,w}\quad(p(w)>0).\]
The zero row uses zero. Thus no extra closure operation is concealed in the formula.

For an outer native map $r$ and admitted continuation maps $s_w$, the pushed joint law has outer law $p'=r_*p$ and row at $w'$ equal to
\[\sum_{w:r(w)=w'}\frac{p(w)}{p'(w')}\,(s_w)_*q_w\]
when $p'(w')>0$. Each summand is in the target continuation specification; finite convexity makes the row admissible. Actual coefficient maps tensor chronologically. This proves bind transport even across noninjective raw maps.

Convexity is material. With two outer witnesses of mass $1/2$ collapsed to one, the continuation specification $\downarrow\delta_0\cup\downarrow\delta_1$ admits the two separate choices but not their merged fair mixture. Downward closure alone therefore fails the transport obligation.

Associativity identifies the same retained triples with laws $p(w)q_w(z)r_{w,z}(t)$; either nesting supplies exactly those choices. The singleton specification $\downarrow\delta_*$ is a unit: smaller unit weights are absorbed by lowerness, and weight one recovers every admitted law. These proofs establish the native chronological construction, not unrestricted independent tensor coherence.
\section{Global writer grades and genuine local products}
Let $\mathcal V$ also be monoidal and let $\mathcal D$ be finitely complete and monoidal. A process $A\to B$ retains a context specification at each $a$, one global writer grade $Q=(O,(D_o))$, and full body $W_a\to B\times O$. Native maps keep the actual context maps and one global content map, with the full body square. Sequential composition uses the bind above, chronological $\mathcal V$ tensor and global writer product $\boxtimes$ formed from $\mathcal D$'s monoidal operation. Associativity and units use the actual monoidal constraints, not quotienting grades.

The product of two processes with the SAME endpoints $A,B$ has global content grade $Q_1\times Q_2$ in $\mathrm{Fam}(\mathcal D)$. At $a$ its raw witnesses are pairs $(w_1,w_2)$ agreeing on returned $b\in B$, its coefficient is $C_{w_1}\times C_{w_2}$, and its laws are exactly the joint sublaws whose two marginals lie in the given specifications. The body returns $b$ and both content labels. The universal property follows by the unique raw pairing, the actual coefficient products and global content product; all law constraints are the two projection conditions. In particular, unused global content labels remain.

The terminal hom process has singleton terminal global grade, raw carrier $B$ at every input, terminal coefficients, all finite measures on $B$, and the evident output body. These local products therefore exist as actual native hom products, not only at the level of marginal behaviors.

As in the existing local-product interaction construction, the projections derive a directed sequential comparison
\[(f_1\times f_2);(g_1\times g_2)\longrightarrow(f_1;g_1)\times(f_2;g_2).\]
Projecting a chronological joint law to one strand forms finite mixtures of allowed continuation rows, exactly as in bind transport. The coefficient arrow is the projection-derived map
\[(C_1\times C_2)\otimes(E_1\times E_2)\to(C_1\otimes E_1)\times(C_2\otimes E_2),\]
and the global content map is its analogous actual product comparison. No product preservation by monoidal tensor is assumed. Coherence follows because both pastings have the same projections, whose chronological equations already hold. The product unit is the terminal hom process, not chronological identity. This supplies the hypotheses for the established local-product interaction mechanism, not a new unproved normal-duoidal identification.
\section{Event restriction and shared assembly}
For every label map $k:X\to Y$, extension $k_!$ composes the raw label map and leaves witnesses, coefficients and lawsets unchanged. For $i:E\hookrightarrow X$, use this extension $i_!$. Restriction uses $W_E=\ell^{-1}(E)$ with inclusion $j$, inherited coefficients, and
\[K_E=\{t:j_*t\in K\}.\]
Lowerness gives $K_E=\{\mu|_{W_E}:\mu\in K\}$. It remains nonempty, lower and convex. Every native map $i_!A\to B$ factors uniquely through $B|E$, with the same law-admissibility and coefficient equations. Thus $i_!\dashv i^*$. Consequently event restriction preserves the actual finite limits of common-labelled native specifications constructed above.

For a commuting square $ki=i'h$ of event inclusions, with ambient map $k:X\to X'$ and restricted map $h:E\to E'$, the canonical comparison
\[h_!i^*\Longrightarrow i'^*k_!\]
includes witnesses satisfying the original event into those satisfying the transported event. Its law map is admissible by lowerness. It is invertible for a pullback event square. Pasting is the same nested raw inclusion with unchanged coefficient maps. For diagonal events in $U\times U$ and ambient map $h\times h$, the extra witnesses are precisely transitions whose endpoints become equal only after state identification. This recovers the geometric source of the earlier closure defect.

The shared-label qualification is essential. Joint assembly here uses one declared event interface. After forgetting the state label, products can pair different state values; independently closing then pairing is not the same operation. Nor do right-adjoint event restriction and its Beck--Chevalley maps establish effectful trace dinaturality for all law specifications. The latter concerns chronological cyclic selection and fails for the unrestricted rowwise closure below.

\paragraph{A dinaturality obstruction.} Let $U=1$, $V=\{0,1\}$, all coefficients and writer grades trivial. Let $f:U\to V$ have raw witnesses $0,1$ returning their own labels, and lawset $K_f=\{(a,b)\ge0:a+b\le1\}$. Let $g:V\to U$ be the deterministic unit row at each input. Define rowwise finite closure by admitting a law on closed witnesses whenever each starting-state row is an allowed restricted row. Then closure of $f;g$ over $U$ still has laws $a+b\le1$. Closure of $g;f$ over $V$ instead has $a\le1,b\le1$: the two starting rows may choose $\delta_0$ and $\delta_1$ separately. It admits $(1,1)$, of total mass two. No native admissibility-preserving map can take that law into the first scalar object, since pushforward preserves total mass. Thus the cyclic witness rotation is not an isomorphism: effectful dinaturality fails. Convexity does not remedy the difference between one shared choice and independent choices indexed by the exposed starting state. A trace extension must retain or constrain that cross-input choice dependence.

\paragraph{Shared-label obstruction.} Two identity loops on a two-element state set have two raw witnesses when jointly assembled over the same loop label and then closed; separately closing and taking the unlabelled product produces four, including off-diagonal pairs. The common-label diagonal gives a directed comparison, not a native isomorphism. The event adjunction remains valid: it never asserted that forgetting the common loop label preserved products.
\section{Independent parallel is a different interface}
For this diagnostic take trivial context coefficients and writer grades. Define independent tensor laws by lower convex closure of product laws. Then full history-dependent bind does not give the forward sequential interchanger above for this tensor. Take a fair-bit process on the left, identity on the right, followed by identity on the left and an arbitrary subprobability bit choice on the right. Joint chronological choice can copy the observed left bit into the right choice, giving the diagonal fair law. The tensor of separately composed strands admits at total mass one only fair-left times a right law, hence cannot contain that diagonal law. The canonical output-preserving shuffle in the forward direction is not admissible.

With these trivial coefficients the reverse comparison exists because separate history choices are particular joint-history choices. The analogous raw shuffle on general coefficient categories additionally requires a middle-four interchange, for example the actual symmetries of symmetric monoidal coefficient categories; it is not supplied by bare monoidality. This is not a tensor pseudofunctor or freely exchangeable independent effects. The same-endpoint product construction instead admits all compatible joint laws and has the projection-derived forward comparison. These distinct operations must not be given the same concurrency interpretation.
\section{A different choice interface: globally scoped kernel specifications}
The previous obstruction diagnoses the scope of choice, not merely the size of the lawset. A second construction retains one admissible set of COMPLETE kernels across all inputs. This is different from a family of independently selectable row specifications and must be declared as a different interface.

Assume finite endpoint objects and symmetric monoidal coefficient categories for this construction. The same native process has one nonempty lower convex set
\[K\subseteq\prod_{a\in A}M_{\mathrm{fin}}(W_a).\]
A native cell retains one actual raw map at each input, actual context coefficient arrows and one global content map, and requires the rowwise pushforward of every whole kernel in $K$ to lie in $K'$. No map is chosen separately for each law. Kernel order is componentwise. Principal kernels embed the subpreserving finite-measure doctrine fully faithfully and the normalized exact doctrine fully faithfully; arbitrary exact finite-measure maps again embed only faithfully.

For a nonempty set $S$ of kernels, let $H(S)$ be its downward finite convex hull. Compose two specifications by
\[K;L=H\{B(p,q):p\in K,q\in L\},\qquad
 B(p,q)_a(w,z)=p_a(w)q_{b(w)}(z).
\]
Crucially, $q$ is ONE complete continuation kernel, shared wherever its input $b$ is revisited. It is not selected independently for each hidden witness $w$. Independent tensor is $H\{p\otimes q\}$. Raw witnesses, coefficients and global grades are those of the original native writer composition and tensor.

\paragraph{Hull calculation.} If $F$ is positive bilinear on the relevant kernels, then
\[H(F(H(S),H(T)))=H(F(S,T)).\]
For the nontrivial inclusion, bound each argument by a finite convex combination, use monotonicity, and expand by bilinearity to a finite convex combination with product coefficients. The reverse inclusion follows from the generators. Bind and tensor satisfy this calculation. Therefore their associativity, units and tensor interchange lift from the individual-kernel equalities using the same actual raw tuple bijections and coefficient/content coherence. Native-map transport follows by pushing each generator; the full body square ensures the continuation kernel is evaluated at the same external input. Downward convex closure then preserves admissibility. The identity is the principal identity kernel. This gives the symmetric monoidal native writer structure with one complete continuation kernel at each composition generator; closure does not silently replace that generator rule by arbitrary separately normalized selections at each hidden history.

\paragraph{Finite feedback.} For $f:A\times U\to B\times U$, restrict a whole kernel to its closed witnesses and collect the retained pairs $(u,w)$ in the row at $a$. Write this positive linear map as $L_U$. Define closure by $L_U(K)$. This image is already lower and convex: if $t\le L_U(p)$, put $t$ on the corresponding closed source coordinates and zero on the others to obtain $p'\le p$ with $L_U(p')=t$. Finite $U$ preserves finite support. Equivalently,
\[L_U(H(S))=H(L_U(S)).\]
The global writer grade and coefficients remain unchanged.

For effectful dinaturality, the kernels of the two closed composites are generated by $L(B(p,q))$ and $L(B(q,p))$ for the SAME global choices $p,q$. Cyclically rotating the retained witness pair identifies these generators by commutativity of nonnegative multiplication and the actual coefficient/content symmetries. The hull calculation therefore gives an isomorphism of the entire admissible law specifications. Tightening, superposing and nested vanishing follow by the same generator calculation. Cups and caps are the principal kernels of the finite-measure construction; the snake and duality coherence cells are unchanged. Thus its coherent finite compact feedback lifts to this globally scoped choice interface.

The rowwise counterexample now gives the simplex on both sides: the one kernel of $f:1\to V$ chosen from its specification cannot be reselected separately for two starting rows after cycling. This is a genuine repaired law, obtained by retaining choice scope, not by denying the earlier counterexample.

\paragraph{Shared products and their boundary.} Same-endpoint hom products use the same output-compatible raw pairs and actual context/global-content products as before. Their admissible COMPLETE joint kernels are precisely those whose two marginal complete kernels belong to the original specifications. Native cone factorization proves the product universal property. Composition again gives the projection-derived directed shared-product comparison. The common-event restriction argument applies to global kernel specifications by retaining the input label as part of the labelled carrier; finite inputs identify a kernel with a finite measure on that disjoint union. It preserves labelled joint limits. Forgetting loop labels still does not turn shared closure into independently closed products: the two-versus-four-witness obstruction remains.

\paragraph{Pure input copying is not product-exact.} Pure precomposition along $h:A'\to A$ takes the lower convex hull $H\{h^*p:p\in K\}$, not just its direct image: duplicated rows need not form a lower set. This action follows also from composition with the principal pure kernel. It does not preserve every native hom product. Take $A=1$, $A'=\{0,1\}$ and $h$ the collapse. Let $f,g$ each have constant external output, a fair internal bit carrier, and their principal fair specifications; content labels may retain each bit. Their product admits all fair joint subcouplings. Precomposing this product makes both copied rows bounded by the same chosen global joint law (or its finite convex mixture). By contrast, the product after precomposition admits the full diagonal coupling in one row and the full anticorrelated coupling in the other. No single mass-at-most-one joint law can dominate both distinct mass-one laws. The canonical product comparison therefore fails to be an isomorphism. This is a source-level failure of the pure-input product-exactness required by the earlier stronger shared-state interaction theorem, not a failure of the local products or their projection-derived sequential comparison.

\paragraph{Semantic cost.} The global constructor does not contain every row-adaptive bind from the preceding section. A continuation can depend on its actual intermediate interface $b$, and histories explicitly exported into that interface remain available. Arbitrary separately normalized continuation selections at hidden histories are not automatically admitted by that generator rule. Downward closure nevertheless permits history-dependent loss masks, and convexification permits mixtures of entire plans. Excess unnormalized mass may make such masking powerful; no blanket prohibition of adaptive behavior is claimed. For a fixed normalized fair source, masking cannot both enforce perfect copying into an independent globally chosen subprobability continuation and preserve total mass one. Such adaptation requires declared interface/state/grade structure or an additional policy semantics. The construction proves a coherent family member supporting finite trace and genuine shared native products; it does not establish unrestricted adaptive concurrency, arbitrary polynomial content, a free completion or infinite-effect laws.

\section{Attribution and remaining interfaces}
The categorical projection-derived interaction mechanism is already developed in
\path{foundations/shared-witness-interaction/Local_Product_Shared_Interaction.tex}.
The contribution under investigation here is a law-specification source of its native products compatible with event restriction, together with the exact failure of the single-law family alternative. No novelty claim is made about finite limits or powersets of laws.

Varacca--Winskel, \emph{Distributing probability over nondeterminism}, distinguishes indexed valuations from geometrically convex alternatives and relates the difference to scheduler power:
\url{https://www.lacl.fr/~dvaracca/papers/distrib.pdf}.
Mio--Vignudelli study the established convex distribution monad in \emph{Monads and Quantitative Equational Theories for Nondeterminism and Probability}:
\url{https://drops.dagstuhl.de/entities/document/10.4230/LIPIcs.CONCUR.2020.28}.
The present native-carrier, coefficient and global-grade proofs are not attributed verbatim to those sources.

A free/universal completion characterization has not been proved. The row-adaptive and globally scoped constructions have distinct choice interfaces; they are not interchangeable interpretations. The writer model is unary content, not arbitrary polynomial dependent occurrences. The row-adaptive construction admits all its declared history-dependent mixtures, whereas the global construction uses the whole-kernel generator rule; no physical scheduler or observation admissibility is inferred. Event restriction is not a proof of recurrence, infinite-effect laws, unrestricted trace or the whole governing v0.2 ambition. The formulation exposes these outstanding interfaces rather than counting local finite completeness as their solution.
\end{document}
